\documentclass[11pt]{article}

\usepackage[margin=1in]{geometry}
\usepackage{booktabs}
\usepackage{hyperref}
\usepackage{amsmath}
\usepackage{enumitem}
\usepackage{xcolor}

\hypersetup{
    colorlinks=true,
    linkcolor=blue,
    urlcolor=blue,
    citecolor=blue
}

\title{INJECT: Science Use Cases and Pipeline Overview}
\author{Sneha Nair}
\date{\today}

\begin{document}
\maketitle

\begin{abstract}
INJECT is a Rubin/LSST-oriented artificial star-cluster injection pipeline for injection-recovery experiments, detector benchmarking, and completeness studies. The package supports smooth analytic cluster profiles, discrete-star cluster generation, Rubin Science Platform workflows, local exploratory workflows, user-defined detection functions, and reproducible output bookkeeping. This document summarizes the science motivation, supported use cases, main pipeline stages, and configuration choices for collaborators and project planning.
\end{abstract}

\section{Project Motivation}

Artificial-source injection is a direct way to measure what a detection or measurement pipeline can recover from real astronomical images. INJECT focuses on synthetic star clusters rather than isolated point sources. The core science question is:

\begin{quote}
Given a population of clusters with known input properties, which objects are recovered by a detector, and how does recovery depend on magnitude, size, profile shape, band, local background, and PSF behavior?
\end{quote}

The package is designed for Rubin-era cluster searches where completeness and failure modes must be measured under realistic image conditions.

\section{Primary Science Use Cases}

\subsection{Detector Benchmarking}

INJECT can be used to compare detection methods on the same injected truth catalog. A user can inject clusters, run a detector, match detections back to injected positions, and compare recovery metrics across detector versions or parameter choices.

Recommended setup:
\begin{itemize}[leftmargin=*]
    \item Start with one band and a small number of injected clusters.
    \item Use fixed random seeds while tuning the detector.
    \item Save both injected truth catalogs and detection catalogs.
    \item Compare completeness versus magnitude and half-light radius.
\end{itemize}

\subsection{Completeness Studies}

Completeness studies ask what fraction of injected clusters are recovered as a function of input properties. The most common completeness metric is

\begin{equation}
    C(x) = \frac{N_{\mathrm{recovered}}(x)}{N_{\mathrm{injected}}(x)},
\end{equation}

where $x$ may be magnitude, half-light radius, concentration, color, background level, or another science variable.

Recommended setup:
\begin{itemize}[leftmargin=*]
    \item Use broad but scientifically motivated parameter ranges.
    \item Run repeated batches to populate each bin with enough injected objects.
    \item Store configuration snapshots with each run.
    \item Report binning choices, matching radius, quality cuts, and uncertainty estimates.
\end{itemize}

\subsection{PSF-Sensitive Validation}

Cluster detectability depends on the local point-spread function. On the Rubin Science Platform, INJECT can use Rubin Butler coadds and Rubin-native PSF objects. The pipeline can also record Rubin PSF-quality mask flags for each injection position.

Recommended setup:
\begin{itemize}[leftmargin=*]
    \item Run on RSP when Rubin-native PSF evaluation is required.
    \item Use \texttt{use\_actual\_psf=True}.
    \item Keep PSF mask flag recording enabled.
    \item Decide whether to keep flagged locations for auditing or skip them with \texttt{skip\_bad\_psf\_regions=True}.
\end{itemize}

\subsection{Multiband Injection-Recovery}

INJECT supports one band, selected band combinations, or all Rubin optical bands. Multiband runs keep injected cluster positions aligned across bands so that recovery can be compared consistently.

Examples:
\begin{verbatim}
InjectionConfig(band="i")
InjectionConfig(bands=["g", "r", "i"])
InjectionConfig(bands=["u", "g", "r", "i", "z", "y"])
\end{verbatim}

Use multiband workflows when detection uses color information, when band-to-band recovery differences matter, or when completeness must be reported across multiple filters.

\subsection{Remote Rubin Access and Local Development}

Not every workflow needs the Rubin stack. INJECT also supports local or TAP-style exploratory runs that use fallback PSF handling. This is useful for development, demonstrations, and early detector experiments before moving to full Butler-backed RSP runs.

\section{Pipeline Overview}

The main pipeline sequence is:

\begin{enumerate}[leftmargin=*]
    \item Configure the run with \texttt{InjectionConfig} and \texttt{ClusterConfig}.
    \item Load image data from an array, TAP-style cutout, or Rubin Butler coadd.
    \item Generate an injected truth catalog with positions and cluster properties.
    \item Render synthetic clusters using smooth profiles or discrete-star generation.
    \item Convolve with Rubin-native or fallback PSF models.
    \item Add injections to the image.
    \item Run a user-supplied detector.
    \item Match detections to the truth catalog.
    \item Compute completeness and recovery summaries.
    \item Save catalogs, checkpoints, and diagnostic plots.
\end{enumerate}

\section{Core Configuration Concepts}

\subsection{Cluster Population}

\texttt{ClusterConfig} controls the injected cluster population:

\begin{itemize}[leftmargin=*]
    \item profile family: \texttt{king}, \texttt{plummer}, \texttt{eff}, or \texttt{sersic}
    \item injection method: \texttt{smooth} or \texttt{discrete}
    \item apparent AB magnitude range
    \item half-light radius range in pixels
    \item concentration range
    \item age range for population bookkeeping and discrete-star workflows
\end{itemize}

\subsection{Run-Level Controls}

\texttt{InjectionConfig} controls run behavior:

\begin{itemize}[leftmargin=*]
    \item run name and output directory
    \item number of clusters
    \item random seed
    \item single-band or multiband selection
    \item cutout size and edge buffer
    \item PSF strategy and fallback FWHM
    \item PSF mask flag recording and bad-region skipping
    \item Butler provenance such as tract and patch
\end{itemize}

\subsection{Magnitude and Distance Modulus}

The pipeline injects sources into observed images, so \texttt{mag\_min} and \texttt{mag\_max} are apparent AB magnitudes. If a science case starts from absolute magnitudes, users should convert first:

\begin{equation}
    m = M + \mu + A_{\mathrm{band}},
\end{equation}

where $m$ is apparent magnitude, $M$ is absolute magnitude, $\mu$ is distance modulus, and $A_{\mathrm{band}}$ is band-specific extinction. The distance modulus is

\begin{equation}
    \mu = 5 \log_{10}\left(\frac{d}{10\,\mathrm{pc}}\right).
\end{equation}

For example, at $d=800{,}000$ pc, $\mu \approx 24.52$. An absolute magnitude range of $M=-6$ to $M=-2$ becomes an apparent magnitude range of approximately $m=18.52$ to $m=22.52$, before extinction.

\section{Execution Modes}

\begin{table}[h]
\centering
\begin{tabular}{lll}
\toprule
Mode & Best For & PSF Behavior \\
\midrule
Local array input & Development and unit tests & Analytic or fallback PSF \\
TAP/local workflow & Lightweight Rubin-adjacent exploration & Metadata-driven fallback PSF \\
RSP Butler workflow & Coadd-based Rubin science runs & Rubin-native CoaddPsf when available \\
Batch workflow & Completeness campaigns & Shared PSF cache and checkpointing \\
Multiband workflow & Color-aware recovery studies & Per-band image and PSF context \\
\bottomrule
\end{tabular}
\caption{Common INJECT execution modes.}
\end{table}

\section{Detector Interface}

The detector is intentionally user-supplied. A detector function should accept an injected image and return a list of detections. Each detection must include at least pixel coordinates:

\begin{verbatim}
def my_detector(image):
    return [
        {"x": 143.2, "y": 287.5},
        {"x": 412.0, "y": 98.4, "snr": 5.4},
    ]
\end{verbatim}

Optional fields such as magnitude, flux, half-light radius, signal-to-noise ratio, ellipticity, or quality flag enable richer recovery diagnostics.

\section{Expected Outputs}

Typical outputs include:

\begin{itemize}[leftmargin=*]
    \item injected truth catalogs
    \item detection catalogs
    \item match and retrieval summaries
    \item completeness curves
    \item two-dimensional completeness maps
    \item per-iteration batch checkpoints
    \item diagnostic plots and postage stamps
\end{itemize}

For reproducibility, each run should preserve the full configuration, random seed strategy, code version, detector version, and output directory layout.

\section{Deployment and Documentation}

The project documentation is built with MkDocs. The Read the Docs configuration is stored in \texttt{.readthedocs.yaml}; it builds \texttt{mkdocs.yml}, installs documentation requirements, installs package requirements, and installs the package itself. The Read the Docs site therefore deploys the MkDocs documentation under \texttt{site\_docs/} plus API pages generated through MkDocstrings.

This Overleaf document is separate from the Read the Docs deployment. It is intended as a collaborator-facing science and methods summary. If a PDF version should be part of the public documentation site, export it from Overleaf and add it to \texttt{site\_docs/} as a static asset, then link it from the MkDocs navigation.

\section{Recommended Reporting Checklist}

For each science run, report:

\begin{itemize}[leftmargin=*]
    \item data source, Butler repo, collection, tract, patch, and band selection
    \item cluster profile family and injection method
    \item apparent magnitude range, and any distance modulus or extinction assumptions
    \item half-light radius and concentration ranges
    \item number of injections and number of repeated iterations
    \item random seed strategy
    \item PSF strategy, PSF cache settings, and PSF mask flag handling
    \item detector function and detector parameters
    \item matching radius and quality cuts
    \item completeness binning and uncertainty estimates
\end{itemize}

\end{document}